\documentclass{article}
\usepackage[T1]{fontenc}
\usepackage{lmodern}
\usepackage{graphicx}
\usepackage{amsmath,amssymb}
\usepackage[a4paper,margin=2cm]{geometry}
\usepackage[absolute,overlay]{textpos}
\setlength{\TPHorizModule}{1mm}
\setlength{\TPVertModule}{1mm}
\usepackage[indonesian]{babel}
\usepackage{hyperref}
\setlength{\emergencystretch}{2em}
% unit: Halaman 36

\begin{document}

% === Halaman 36 (llm) ===

\section*{Program Transposition Cipher dalam Bahasa Python (26 huruf alfabet)}

\begin{verbatim}
import math

def transposition_cipher_encrypt(plaintext, k):
    ciphertext = [''] * k    # membuat sebuah list dengan panjang k, yang setiap
                               # elemennya berupa string kosong.

    for kolom in range(k):
        index = kolom
        while index < len(plaintext):
            ciphertext[kolom] = ciphertext[kolom] + plaintext[index]
            index = index + k
    return ''.join(ciphertext)
\end{verbatim}

\end{document}
